% FRONTEND_DESIGN.tex -- what the LBM Solver Studio frontend must carry for the
% WB HOME4 D3Q27 hydro + phase-field stack, and what would make it exceptional.
% Design note, 2026-10-05.  Compile: latexmk -pdf FRONTEND_DESIGN.tex
\documentclass[11pt,a4paper]{article}
\usepackage[margin=2.2cm]{geometry}
\usepackage{amsmath,amssymb,bm}
\usepackage{booktabs,longtable,array,multirow}
\usepackage{siunitx}
\usepackage[expansion=false]{microtype}
\usepackage{enumitem}
\usepackage[colorlinks=true,linkcolor=blue!50!black,citecolor=blue!50!black,urlcolor=blue!50!black]{hyperref}
\usepackage{xcolor}
\usepackage{float}
\newcolumntype{L}[1]{>{\raggedright\arraybackslash}p{#1}}
\usepackage{listings}
\setlist{nosep,leftmargin=1.5em}
\sisetup{range-phrase=--,range-units=single,per-mode=symbol}
\lstset{basicstyle=\ttfamily\small,columns=fullflexible,keepspaces=true,breaklines=true,frame=single,framerule=0.3pt}
\newcommand{\Ma}{\mathrm{Ma}}
\renewcommand{\Re}{\mathrm{Re}}
\newcommand{\Fr}{\mathrm{Fr}}
\newcommand{\Bo}{\mathrm{Bo}}
\newcommand{\We}{\mathrm{We}}
\newcommand{\Ca}{\mathrm{Ca}}
\newcommand{\Pe}{\mathrm{Pe}}
\newcommand{\Cn}{\mathrm{Cn}}
\newcommand{\Kn}{\mathrm{Kn}}
\newcommand{\At}{\mathrm{At}}
\newcommand{\KC}{\mathrm{KC}}
\newcommand{\Ga}{\mathrm{Ga}}
\newcommand{\cs}{c_s}
\newcommand{\ps}{p^\ast}
\newcommand{\pd}{p_d}
\newcommand{\Pih}{\Pi_h}
\newcommand{\Mcells}{\,\mathrm{Mcells}}
\newcommand{\code}[1]{\texttt{#1}}
\newcommand{\todo}[1]{\textcolor{red!70!black}{[#1]}}
\newcommand{\have}{\textcolor{green!45!black}{exists}}
\newcommand{\addme}{\textcolor{orange!80!black}{add}}
\newcommand{\gui}{\textcolor{blue!55!black}{GUI}}

\title{\textbf{FRONTEND\_DESIGN} --- what \emph{LBM Solver Studio} must carry for the
WB HOME4 D3Q27 hydro + phase-field stack, and what would set it apart\\[2pt]
\large Design note for the frontend build}
\author{HOME4 project note}
\date{2026-10-05}

\begin{document}
\maketitle

\begin{abstract}
\noindent The first prototype of \emph{LBM Solver Studio} is a clean, well-built shell
whose \emph{shape} was borrowed from finite-volume Navier--Stokes products (Fluent,
SU2, STAR): a meshing step, material properties, a solver page, residual-style
monitors, and a post-processor. HOME4 has none of the first three in that form and
offers several things an FV code cannot. This note gives the frontend engineer the
dozen facts about the solver that should shape the interface, a specific critique of
the prototype, a replacement information architecture, and inventories --- fields,
monitors, forces, zones, multidomain, run/backends, outputs, viewer --- each mapped to
the solver attribute or driver flag that supplies it, marked \have\ (the code
already produces it) or \addme\ (a small solver-side addition is needed). It closes
with the ranked list of features that would make the product different from every
other CFD front end, a phased build sequence, a proposed telemetry schema so the GUI
never has to import \code{torch}, and a worked map of the \path{run_hull_speed.py}
flags onto controls. The thesis in one sentence: \textbf{organise the interface
around lattice units, ledgers and budgets, the diffuse interface, and validated
recipes --- not around meshing, residuals and materials.}
\end{abstract}

\tableofcontents

%======================================================================
\section{Scope}
%======================================================================
This note covers the 3-D stack in \path{home4_hydro_phase_3d}: the single-domain
velocity-based D3Q27 solver (\path{home4_3d_multiphase.py}), the well-balanced
layer (\path{home4_3d_wb.py}), the FSI and moving-body layers
(\path{home4_3d_fsi_solver.py}, \path{home4_3d_moving_fsi.py},
\path{hull_moving.py}), the multidomain stack (\path{home4_3d_md_*.py},
\path{multidomain_nlevel.py}), the CAD$\to$SDF pipeline (\path{cad_sdf.py}), the
run drivers (\path{run_*.py}), the PyVista viewer (\path{kinetic_viz/}) and the
rack workflow (\code{rack/}, \path{RACK_HANDOFF.md}). The mathematical content is
in \path{home4_3d_hydro_phase_math.tex}; equation references below point there.
The 2-D stack is out of scope but the same architecture applies.

The note is written for someone who will build the interface without having read
the solver. Where it says ``the solver stores $X$'' it names the attribute. Where it
proposes a new diagnostic it says what has to be added and roughly where.

%======================================================================
\section{Twelve facts about the solver that should shape the interface}\label{sec:facts}
%======================================================================
\begin{enumerate}[leftmargin=2em]
\item \textbf{Everything runs in lattice units.} Lengths in cells, time in steps,
  speed in cells/step with $\cs = 1/\sqrt3$. Physical units are a \emph{map} chosen at
  setup ($\Delta x$, $\Delta t$, $\rho_{\rm ref}$), not a property of the fields. The
  user, however, thinks in $\Re$, $\Fr$, $\Bo$, $\We$, $\Ca$, $\Pe$, $\At$, $\KC$,
  $\Ga$; every benchmark is specified and reported that way (math note \S2.2).
\item \textbf{There is no mesh.} There is a Cartesian lattice ($N_x\times N_y\times N_z$),
  a signed-distance band around CAD bodies (\path{cad_sdf.py}: tessellate, weld,
  orient, band distance, winding-number sign, EDT completion), cut links with
  interpolation weights, and multidomain refinement patches with factor-2 levels.
  ``Meshing'' as a stage does not exist; ``Lattice'' and ``Geometry'' do.
\item \textbf{The scheme is explicit and time-accurate.} There are no iterations
  within a step, hence no residual to converge. The natural health signals are
  conservation ledgers, budget closures and agreement between independent channels
  (\S\ref{sec:monitors}).
\item \textbf{The zeroth moment is pressure, not density.} The hydro population $g$
  carries $\ps$ (normalised pressure); $\rho$ is a function of $\phi$
  (\path{update_density_from_phi}). Under the WB layer $\ps$ \emph{is} the dynamic
  pressure $\pd$ and the hydrostatic head $\Pih$ is carried separately; physical
  pressure is reassembled as $p = \rho\cs^2\pd + (\Pih-\Pih^0)$
  (\path{WBFSI3D.measure_forces}).
\item \textbf{The strain rate is a state variable.} $S_{ab}$ (six components) is
  stored and relaxed in closed form; $u$, $\ps$, $S$ and the seven tied third-order
  closures $a^{(3)}$ are the hydro state. Dissipation $2\nu S\!:\!S$, enstrophy and
  $Q$ come without finite differencing.
\item \textbf{Relaxation time is a field.} $\tau(\phi)$ (\path{tau_fld}) runs
  between $\tau_H$ and $\tau_L$; the light phase is the dangerous end
  ($\nu_L$ small, $\tau_L\to\tfrac12$).
\item \textbf{The interface is a diffuse object with its own parameters.} $\xi$ in
  cells (4--5), $\sigma$, mobility $M$, four phase relaxation rates
  $(s_T, s_{D1}, s_{D2}, s_{xy})$, the \code{muphi} force form, and two
  load-bearing safeguards (gradient limiter at $1.6\times$ the tanh slope; force
  thresholding at $60|F_b|$ with factor $0.6$ where $\phi<0.1$; math note \S5.3).
\item \textbf{Forces are a decomposition.} $F = F_s + F_p + F_\nu + F_b$
  (single domain) or $F_s + F_p(\pd) + F_\nu + F_{\rm tilt}$ (WB), each a
  $[3,N_x,N_y,N_z]$ field the user may want to see separately.
\item \textbf{Bodies are SDF bands, and moving them has a cost cliff.} Static mask:
  47.8 MLUPS; per-step retabulation: 5.9 (oscillating cylinder) and 0.63
  (sedimentation) on the M1 Pro (math note, performance summary table). Two
  independent force channels exist: stress integral and momentum exchange.
\item \textbf{Multidomain has per-level physics.} With $k=2^d$:
  $\nu^{(d)}=\nu_0 k$ (or a $\tau$ floor), $\sigma^{(d)}=\sigma_0 k$,
  $M^{(d)}=M_0 k$, $g^{(d)}=g_0/k$; $\xi$ is fixed in \emph{root} cells; mobility is a
  fine-level quantity; each level keeps a mass ledger (math note \S10.5, Prop.~10.1).
\item \textbf{Three backends, one of them silently slow.} Metal (Mac), CUDA
  (rack), or the pure-PyTorch fallback when the extension fails to import --- tens of
  times slower with no other symptom than a stdout line. The code is
  bandwidth-bound ($\sim$500 B/node/step for two-buffer D3Q27).
\item \textbf{Four kinds of output already exist.} Traces (every $N$ steps),
  mid-span slices, full 3-D viz snapshots (\path{kinetic_viz.data.save_snapshot}:
  \code{ux}, \code{uy}, \code{uz}, \code{phi}, \code{solid}, \code{iteration}, \code{spacing}, \code{origin}; atomic rename), and restart
  checkpoints (\path{save_state}: the full moment state incl.\ $S$, $a^{(3)}$,
  $J_\phi$, $P_\phi$). They are different things and belong on the timeline as
  different markers.
\end{enumerate}

\begin{table}[H]
\centering\small
\caption{Translation table for the frontend engineer: the FV concept the prototype
was modelled on, and the HOME4 counterpart.}
\begin{tabular}{L{4.2cm}L{10.8cm}}
\toprule
FV-GUI concept & HOME4 counterpart \\
\midrule
Mesh / meshing & Lattice resolution + unit map; SDF band from CAD; cut links; MD levels \\
Material properties & Two fluids ($\rho_H,\rho_L,\nu_H,\nu_L$) and an interface ($\sigma,\xi,M,s^\phi$) \\
Residual convergence & Ledgers (mass, energy), channel agreement (forces), rest tests (WB) \\
CFL number & Lattice Mach $U/\cs$ and $\tau$ margin; interface $\Pe$ and $\Cn$ \\
Pressure field & $\ps$, $\pd$, $\Pih$, reassembled $p$ --- four layers, not one \\
Velocity gradient (post-processed) & $S_{ab}$ stored; dissipation/enstrophy/$Q$ exact from state \\
Turbulence model panel & None; instead the kinetic layer ($S^{\rm neq}$, $a^{(3)}$, filtered $a^{(4)}$) \\
Inlet/outlet/wall BCs & Plus sponges, beaches, wave absorption, phase-wall pinning, mass correction \\
Time step size & Fixed by the lattice; the knobs are $N$, $U$ (Ma) and $\tau$ \\
Checkpoint & Viz snapshot \emph{or} restart state; not the same file \\
\bottomrule
\end{tabular}
\end{table}

%======================================================================
\section{The prototype: keep, change, remove}\label{sec:proto}
%======================================================================
\paragraph{Keep.} The visual language (dark, quiet, dockable panels with
collapse buttons), the top bar (Run / Pause / Stop / Step / Export), the playback bar
with a snapshot index, the axes gizmo, the colour-scale panel with a ``source
range'' line, the view presets (Flow overview, Perspective), the notification
bell, and the recovery banner. \emph{Step} is especially right for an LBM: one step
is cheap and well-defined; add ``step $N$'' and ``run to $t^\ast=\ldots$''.

\paragraph{Change.}
\begin{itemize}
\item The colour bar reads \qtyrange{0.005}{397}{m/s} for an airfoil. At lattice
  Mach above $\sim$0.3 this solver diverges; 397 m/s is supersonic in air. The demo
  numbers show the shell has no notion of the unit map. Every scalar readout needs
  lattice, physical and nondimensional forms, with one global toggle and the other
  two in a tooltip.
\item The playback bar shows ``frame 600, 0.1200 s''. Show step, physical time, and
  nondimensional time ($t/T$, $tU/L$, $t/t_0$) together; the last is the one every
  validation plot is drawn against.
\item The demo project is ``Airfoil SU2 009'' (single-phase, compressible, FV). The
  shipped example projects should be the validated recipes (\S\ref{sec:recipes}).
\item ``Meshing'' $\to$ \emph{Lattice}. ``Materials'' $\to$ \emph{Fluids \& Interface}.
  ``Boundary Conditions'' $\to$ \emph{Boundaries \& Zones}. ``Solver'' $\to$
  \emph{Run}. ``Results'' and ``Post-Processing'' merge into \emph{Fields} and
  \emph{Reports}; add \emph{Bodies} and \emph{Validation}.
\item The status line ``RECORDED CFD'' is honest and should stay as a mode badge:
  \emph{live}, \emph{replay}, \emph{remote} (job on the rack).
\end{itemize}

\paragraph{Remove.} Nothing visual. Remove the assumption that a project is one
geometry + one mesh + one solver run; a HOME4 project is a \emph{recipe lineage}
(\S\ref{sec:recipes}) with many runs, each a provenance \path{.npz}.

%======================================================================
\section{Information architecture}\label{sec:ia}
%======================================================================
Sidebar, top to bottom, and what each page owns:
\begin{description}[style=nextline]
\item[Dashboard] Open projects, running jobs (local and remote), last ledgers, last
  failures with their locator (\S\ref{sec:locator}).
\item[Projects] Recipe lineage tree; each run as a card with its \path{.npz}, tag,
  backend, MLUPS, and gate status.
\item[Geometry] CAD import (IGES/STEP/STL/OBJ via \path{cad_sdf.py}), attitude
  (heel/trim/yaw/sinkage), flotation solve (\path{--float}), SDF band preview against
  the voxel staircase, cut-link fraction, tessellation cache status.
\item[Lattice] $N_x,N_y,N_z$; unit map; pads/air gap/depth in body lengths
  (\path{--pad_up}, \path{--pad_down}, \path{--pad_side}, \path{--depth},
  \path{--air}); MD levels with per-level parameters; memory and cost estimate.
\item[Fluids \& Interface] $\rho_H,\rho_L,\nu_H,\nu_L,\sigma,\xi,M$,
  $s^\phi$, \path{tau_method}, \path{surface_tension_form}, safeguards; the
  dimensionless-first form (\S\ref{sec:units}) lives here.
\item[Boundaries \& Zones] Walls, inlet/outlet with ramp, periodic axes, sponges and
  beaches, wave absorption, phase-wall pinning ($\phi_{\rm top},\phi_{\rm bot}$),
  mass-correction mode, drawn as translucent regions in the 3-D view.
\item[Bodies] Fixed / forced (heave, pitch, spin) / free; mass properties, CoG,
  stiffness matrix $K$; retabulation policy and its cost; force channels.
\item[Run] Backend, device, block shape, output cadence (\path{--measure_every},
  \path{--print_every}, \path{--viz-every}, \path{--save_state}), warm start
  (\path{--init_state}), tag, outdir, queue target (local / gpulab), live CLI
  equivalent.
\item[Monitors] Ledgers and budgets first; forces; performance; safeguard activity.
\item[Fields] The 3-D viewer (\S\ref{sec:viewer}).
\item[Validation] Recipe cards with reference overlays and gates
  (\S\ref{sec:recipes}); convergence-ladder builder.
\item[Reports] Figures with sidecars, tables, the run's CLI line and parameter
  block; export to the LaTeX report conventions.
\item[Settings] Unit display, palette, viewer defaults.
\end{description}

%======================================================================
\section{Setup: dimensionless first, with a live feasibility map}\label{sec:units}
%======================================================================
The user enters the dimensionless groups of the case and the lattice resolution;
the GUI derives the lattice parameters and checks them. The relations are those
already used in the drivers (\path{run_hull_speed.py} prints exactly this block at
startup: $U$ in lu/step, $\tau$, $\Re_L$, $U/\cs$, $\Fr_h$):
\begin{align}
\nu &= \cs^2\,(\tau-\tfrac12), \qquad \Ma = U/\cs, \qquad \Re = \frac{UL}{\nu}
      = \frac{\Ma\,L}{\cs(\tau-\tfrac12)}, \\
\Kn &\approx \frac{\Ma}{\Re} = \frac{\cs(\tau-\tfrac12)}{L}, \qquad
\Cn = \frac{\xi}{L}, \qquad \Pe_\phi = \frac{UL}{M}, \\
\Fr &= \frac{U}{\sqrt{gL}}, \qquad \Bo = \frac{\rho_H g L^2}{\sigma}, \qquad
\We = \frac{\rho_H U^2 L}{\sigma}, \qquad \At = \frac{\rho_H-\rho_L}{\rho_H+\rho_L}.
\end{align}
For fixed $\Ma$ and $\tau$, $\Re\propto L$: resolution buys Reynolds number. The
three-way trade ($N$, $\Ma$, $\tau$) is the central fact of LBM sizing and should be
a \emph{live widget}: drag any one, the others respond, and the validated recipes
are plotted as points so the safe region is visible. The widget replaces the
Knudsen number as a display: $\Kn$ is derived from two things the user already
controls, is tiny by construction in the hydrodynamic regime the Chapman--Enskog
analysis assumes (math note \S6), and is best shown as a tooltip.

\begin{longtable}{L{3.6cm}L{4.2cm}L{7.2cm}}
\caption{Setup checks the GUI should run before launch. All derive from
parameters already in the drivers.}\\
\toprule
Check & Rule & Message / action \\
\midrule
\endfirsthead
\toprule
Check & Rule & Message / action \\
\midrule
\endhead
Lattice Mach & $U/\cs \le 0.1$ (warn), $\le 0.3$ (block) & ``Compressibility error $\sim\Ma^2$''; suggest raising $N$ or lowering $U$ \\
$\tau$ margin & $\tau_L - \tfrac12 \ge 0.01$ (warn below) & Light phase under-resolved viscously; suggest \path{tau_method}, $\nu_L$, or MD $\tau$ floor \\
Interface width & $4 \le \xi \le 5$ cells & Below: spurious currents and limiter firing; above: $\Cn$ too large \\
Cahn number & $\Cn=\xi/L$ vs the recipe's & Fixed-$\Cn$ refinement means $\xi\propto N$; warn when a ladder holds $\xi$ in cells \\
Mobility & $M$ matches recipe at the \emph{finest} level & MD: pass $M_0 = M_{\rm ref}/2^{d_{\max}}$ (defect D23) \\
Density ratio & $\rho_H/\rho_L \ge 100$ flags safeguards & Turn on gradient limiter and force thresholding (both required at 1000:1) \\
Memory & $\approx 350$ B/node $\times N_{\rm cell}$ (+ MD, FSI) & Compute from the allocation list, not a constant; M4 Pro MPS ceiling $79.5\Mcells$ \\
Index width & Metal path: $27 N_{\rm cell} < 2^{32}$ ($N_{\rm cell} < 159$ M) & Block on Metal; CUDA path uses \path{size_t} \\
Backend & extension imported? & Big badge: Metal / CUDA / \emph{PyTorch fallback (slow)} \\
Sponge sizing & width vs expected wavelength $2\pi\Fr^2 L$ & Hull driver already compares the wake wavelength to $2\pi\Fr^2L$; use the same number to size \path{--sponge} and \path{--beach_gap} \\
\bottomrule
\end{longtable}

%======================================================================
\section{Field inventory}\label{sec:fields}
%======================================================================
Fields the viewer and the slice tools should offer, with the solver attribute that
supplies them. Default air masking follows the lesson that $\omega c/U$ normalisation
amplifies air-phase noise by $\sim$1600$\times$.

\begin{longtable}{L{3.3cm}L{4.6cm}L{1.4cm}L{5.4cm}}
\caption{Field inventory. ``mask'' = air phase masked by default.}\\
\toprule
Field & Source & Status & Notes \\
\midrule
\endfirsthead
\toprule
Field & Source & Status & Notes \\
\midrule
\endhead
$\phi$ & \code{phi} & \have & Isosurface at 0.5; thickness histogram (\S\ref{sec:monitors}) \\
$\rho(\phi)$ & \code{rho} & \have & Linear in $\phi$ \\
$\tau(\phi)$ & \path{tau_fld} & \have & Show $\tau-\tfrac12$ log scale; the light phase is the risk \\
$\ps$ & \path{p_star} & \have & Under WB this \emph{is} $\pd$ \\
$\pd$, $\Pih$, $p$ & \path{p_star}, \path{compute_Pi_h}, \path{measure_forces} & \have & Three-layer pressure view; rest test $\pd=0$ \\
$u$, $|u|$ & \code{ux,uy,uz}, \path{get_velocity_magnitude} & \have & mask for vorticity-derived quantities \\
$S_{ab}$ & \path{Sxx..Syz} & \have & Stored state; no differencing \\
Dissipation $2\nu S\!:\!S$ & from $S$, \path{tau_fld} & \addme & One kernel-free tensor expression; per-phase sums for the budget \\
$\nabla\cdot u$ & \path{compute_vel_gradients} trace & \addme & Compressibility map; the observable for \path{s_bulk} \\
$\omega$, $Q$, helicity & \path{kinetic_viz} (central differences) & \have & Could instead use $S$ + antisymmetric part from the gradient kernel \\
$a^{(3)}$ (7), filtered $a^{(4)}$ (6) & \path{a3_*}, \path{_compute_a4_all} & \have & Kinetic layer; magnitude maps \\
$J_\phi$, $P_\phi$ & \path{Jx_phi..Pyz_phi} & \have & Phase-moment layer; rarely needed, keep behind ``advanced'' \\
$F_s, F_p, F_\nu, F_b, F_{\rm tilt}$ & \path{compute_all_forces} (returns the sum) & \addme & Return the terms separately behind a flag; mask \\
$\nabla\phi$, $|\nabla\phi|$ vs $4\phi(1-\phi)/\xi$ & \path{compute_grad_phi} & \have & Limiter-engagement map = where the ratio exceeds 1.6 \\
Solid mask, SDF, cut links & \path{fsi_solid}, SDF bundle, \code{faces} & \have & Draw the SDF zero level \emph{and} the voxel contour \\
MD level map, bands, ghosts & \code{GridLevel} topology & \have & Colour by level; show band depth and twin rings \\
Sponge / beach weights & \path{_build_absorb_field}, driver zones & \have & Translucent volumes with strength profile \\
Safeguard activity & limiter/threshold masks & \addme & Boolean per cell per step; count + map \\
\bottomrule
\end{longtable}

%======================================================================
\section{Monitors: ledgers and budgets before anything else}\label{sec:monitors}
%======================================================================
There is no residual. The Monitors page should open on five health signals, each
with a traffic-light threshold and a one-line explanation of what to do when it
turns.

\begin{longtable}{L{3.4cm}L{4.6cm}L{1.4cm}L{5.3cm}}
\caption{Health signals and secondary monitors.}\\
\toprule
Monitor & Source & Status & Healthy / action \\
\midrule
\endfirsthead
\toprule
Monitor & Source & Status & Healthy / action \\
\midrule
\endhead
\multicolumn{4}{l}{\emph{Primary (the five)}}\\
$\phi$ mass ledger & \path{get_phase_mass}; hull driver's ``ledger drift''; MD per-level ledger + injections & \have & $|{\rm drift}| < 10^{-4}$; injections should not grow (positive-feedback failure, Prop.~10.1) \\
Energy budget residual & live version of \path{h2_energy_budget.py}: $W = D_{\rm near}+D_{\rm far}+D_{\rm air}+Z_{\rm beach}+Z_{\rm floor}+\Delta KE+\Delta PE+{\rm res}$ & \addme & Residual as a bar; large residual = non-conservative operations at the body or sub-grid Stokes layer \\
Force channel agreement & \path{measure_forces}: stress integral vs \path{MEA_ALD} & \have & Within a few \%; divergence means cut-link / retabulation trouble \\
Steady-window bracket & driver: window mean [previous window] & \have & Bracket within tolerance = converged in time \\
WB rest test & $\max|\pd|$ with $U=0$, full $g$ & \have & Exactly 0 at rest; any growth is a BC or MD defect \\
\midrule
\multicolumn{4}{l}{\emph{Secondary}}\\
$\Ma_{\rm inst}$, $\max|u|$ + location & \path{get_max_speed}; hull driver \path{vw_loc} & \have & Location feeds the locator (\S\ref{sec:locator}) \\
$\tau_{\min}$ & \path{tau_fld.min()} & \have & Margin to $\tfrac12$ \\
$\|\nabla\cdot u\|$ & new & \addme & $\sim\Ma^2$; rises with acoustic ringing; knob \path{s_bulk} \\
Interface thickness histogram & $|\nabla\phi|$ at $\phi\in(0.1,0.9)$ & \addme & Peak at $\xi$; a tail = sharpening/limiter fight \\
Spurious current & $\max|u|$ in cells far from forcing (Laplace test) & \addme & Recipe values known from the droplet test \\
Safeguard counts & limiter/threshold cells per step & \addme & Non-zero but flat is tolerable; climbing is a warning \\
KE per phase, surface energy & $\tfrac12\rho|u|^2$ masked; $\sigma\times$ area of $\phi=0.5$ & \have/\addme & Feed the budget \\
Non-finite guard & driver ``non-finite load'' check & \have & Auto-stop, auto-checkpoint last good state \\
MLUPS inst.\ (windowed) and cum. & drivers (cum.); windowed to add & \addme & Cum.\ masks throttling; show both \\
Achieved GB/s & bytes/node/step $\times$ MLUPS vs device peak & \addme & The honest efficiency number for a bandwidth-bound code \\
\bottomrule
\end{longtable}

%======================================================================
\section{Forces, bodies, attitude}\label{sec:forces}
%======================================================================
\path{WBFSI3D.measure_forces} returns \code{Cd, Cl, Fx, Fy, Fz} from the stress
integral on the cut faces (with $p = \rho\cs^2\pd + (\Pih-\Pih^0)$), and
\path{MEA_ALD} from momentum exchange when post-streaming populations are kept.
\path{hull_moving.measure_forces_moments} adds the moment about the CoG.
The GUI should show, per body:
\begin{itemize}
\item the time series of $F_x,F_y,F_z,M_y$ with the pressure and viscous parts, in the
  benchmark normalisation ($F/mg$, $M_y/mgL$, $C_d=F_x/(\tfrac12\rho U^2 D\,N_y)$);
\item both channels and their difference;
\item the steady window with the previous window in brackets, and the averaging
  length in body lengths (\path{--avg_L});
\item for floating bodies: the hydrostatic stiffness $K$ ($k_{33},k_{35},k_{55}$) from
  \path{hull_hydrostatics.py}, the held attitude $\zeta_{\rm eq}$, trim, and the
  quasi-static running attitude $K^{-1}F$ next to the tank reference
  (\path{th01_tank_data.py});
\item for forced motion: the fitted added mass and damping (\path{period_fit.py},
  \path{barge_roll_fit.py}) against the BEM reference
  (\path{th01_bem_capytaine.py}) and Vugts;
\item for free motion: the 6-DOF state, virtual-mass integrator status, and the
  retabulation cadence with its MLUPS cost shown next to it.
\end{itemize}

%======================================================================
\section{Boundaries and zones}\label{sec:zones}
%======================================================================
Half the tuning in the hull campaign is zones. Each should be a drawn region with a
strength profile, and the GUI should size them from the expected wavelength:
\path{--sponge} (solver sponge relaxing $\pd$ and $u$ toward the stream),
\path{--xbeach} / \path{--xbeach_strength} (driver beach: $u\to$ stream,
$\phi\to$ flat, $p$ untouched), \path{--beach_y}, \path{--beach_gap},
\path{--zone_strength}, \path{--floor_fric}, \path{--mass_correct}
\{beach, global, off\}, \path{--pierce_bc}, phase-wall pinning
(\path{phi_top}, \path{phi_bot}, \path{pin_phase_walls}), inlet ramp
(\path{--ramp_L}, cubic) and the hull-frame body force $\rho\,dU/dt$
(\path{--no_frame_accel} to drop it). In MD the sponge strength is
exponent-rescaled per level (D13); show the per-level value.

%======================================================================
\section{Multidomain}\label{sec:md}
%======================================================================
A dedicated view: level map coloured by level; per-level table of
$\nu^{(d)},\sigma^{(d)},M^{(d)},g^{(d)},\tau^{(d)}$ with the $2^d$ scaling shown
explicitly and the $\tau$-floor choice flagged as a \emph{physics} choice (it runs
level $d$ at $2^d\times\Re$); band depth, overlap, restriction margin, even-wrap on
periodic axes, phase-free root on/off; the hopping patch following a moving body
(Tier 2); per-level MLUPS and mass ledger; $S^{\rm neq}$ factor mode
(\code{dorschner} / \code{derived}). Driver flags: \path{--levels}, \path{--z1},
\path{--z2}, \path{--margin}, \path{--even_wrap}, \path{--root_phase_free},
\path{--sneq_mode}, \path{--no_tau_floor}, \path{--mass_fix_every},
\path{--no_full_f2c}.

%======================================================================
\section{Run: backends, cost, queue, reproducibility}\label{sec:run}
%======================================================================
\paragraph{Backend badge.} Metal / CUDA / PyTorch fallback, from the extension
import line. Never let a fallback run start without a confirmation.

\paragraph{Cost estimate before launch.} From the MLUPS table (M1 Pro: RTI 84.6,
Colagrossi 32.5, breaking wave $\sim$45, MD RTI 55.3, Magnus 47.8; per-step
retabulation 5.9 / 0.63; rack $\sim$1.0 GLUPS on the ported kernels) and the step
count: ``$\approx$ 4.2 h on the M4 Pro, 35 min on gpulab''. Keep a small per-machine
MLUPS table that the GUI updates from completed runs.

\paragraph{Queue.} Mirror \path{rack/status.sh}: GPU utilisation, memory, power,
temperature, clocks; ``card busy with someone else's job'' detection (the rack is
shared); running jobs with elapsed time; queued jobs; last log lines; disk free. Jobs
are launched as the existing drivers with a generated command line; the GUI tails
the log and parses the telemetry (\S\ref{sec:telemetry}). Local MPS and remote
CUDA from the same dialog.

\paragraph{Reproducibility.} A live ``equivalent command line'' pane that updates as
settings change (the ParaView-trace idea applied to run configuration); the
provenance \path{.npz} as the project record; the \path{_viz} tag rule (figure
runs never overwrite published results, and are checked against the published
coefficients before use); warm start (\path{--init_state}) with the grid check
\path{load_state} already enforces; a JSON sidecar for every exported image
(\path{kinetic_viz} already writes one).

%======================================================================
\section{Outputs and the timeline}\label{sec:outputs}
%======================================================================
Four kinds, four markers:
\begin{enumerate}
\item \textbf{Traces} (\path{--measure_every}): forces, ledgers, $\max|u|$ and
  location, energies; cheap; feed Monitors.
\item \textbf{Slices} (mid-span, as most result archives store): cheap 2-D fields
  at \path{--save_every}; feed the slice tools.
\item \textbf{Viz snapshots} (\path{--viz-every}, \path{--viz-dir}): full 3-D
  \code{ux,uy,uz,phi,solid}; expensive (CPU transfer + compression); feed the viewer;
  one \path{.npz} per snapshot, atomic rename, so the viewer can watch a directory
  safely.
\item \textbf{Restart checkpoints} (\path{--save_state}): the full moment state
  (\path{STATE_KEYS} + $a^{(3)}$); large; ``warm start from here'' on right-click.
\end{enumerate}
The GUI should let the user set cadences for each independently, show the disk cost,
and never conflate a viz snapshot with a restart point. Auto-checkpoint the last
good state when the non-finite guard trips.

%======================================================================
\section{Viewer: adopt the \code{kinetic\_viz} contract}\label{sec:viewer}
%======================================================================
Do not reinvent the viewer; port the layer set and the conventions:
\begin{itemize}
\item Layers: \code{surface} ($\phi=0.5$ isosurface, adjustable), \code{obstacle}
  (mask contour, and the SDF zero level when available), \code{streamlines} (RK45
  tubes, speed-coloured, deterministic seeds in liquid cells), \code{speed} (planar
  slice), \code{vorticity} (RGBA volume: $Q$ colour, $|\omega|^2$ opacity), \code{q}
  (positive-$Q$ isosurface), \code{helicity} (signed plane).
\item Colour ranges fixed to the 99th percentile of the first frame and held across
  frames; explicit \path{--speed-max} / \path{--q-max} for cross-run comparison;
  clipped extremes recorded in the sidecar.
\item Air mask \emph{on by default} for vorticity, $Q$ and pressure panels; a toggle.
\item Derivatives exclude stencils touching air or solid (one-node halo); say so in the
  panel rather than hiding it.
\item Crop in lattice indices with preserved origin; stride as a preview only.
\item MD patches are not auto-stitched; show them with true origin and spacing.
\item Export: PNG + JSON sidecar, VTI for ParaView, PNG sequence for FFmpeg.
\end{itemize}
The frontend's rendering stack is the colleague's choice; the \emph{contract}
(input shape, origin/spacing, masks, percentile ranges, sidecar) is not.

%======================================================================
\section{Validation cards and recipes}\label{sec:recipes}
%======================================================================
A new project starts from a validated recipe card, not a blank canvas. Each card
carries the literature anchor, the parameter recipe (math note \S14), the driver,
the reference data file, and a gate. Users perturb from there; the project keeps the
lineage (``derived from Colagrossi WB; $\Fr$ 0.30 $\to$ 0.45'') and the GUI flags
departures from the validated envelope.

\begin{longtable}{L{2.4cm}L{4.8cm}L{4.0cm}L{3.5cm}}
\caption{Recipe cards to ship.}\\
\toprule
Card & Recipe (from the math note) & Driver / reference & Gate \\
\midrule
\endfirsthead
\toprule
Card & Recipe & Driver / reference & Gate \\
\midrule
\endhead
RTI (Fakhari) & $\At=0.5$, $\Re=3000$, $\Ca=0.26$, $\Pe=1000$, $\xi=5$, $t_0=16000$, $s^\phi=(1,1.5,1.5,1.5)$, $\rho_H{:}\rho_L=3{:}1$ & \path{run_rti3d_multiphase.py} & Spike/bubble fronts vs reference; fixed-$\Cn$ ladder \\
Breaking wave (Banari grid3) & $N_x=400$, $W/L=0.25$, $L_z/L=0.30$, $c_{\rm lat}=0.015$, $\xi=5$, $M=0.02$, slope 0.08 & \path{run_breaking_wave3d_md.py} & Jet geometry and energy vs Banari; $c_{\rm lat}=0.02$ is a known failure \\
Colagrossi WB cylinder & $\Fr$-aware $U$, $g=(U/\Fr)^2/D$, $M=0.05$ ($\Fr<0.8$), $\Bo=200$, sponge $8D$, $s^\phi_{xy}=1.8$, limiter on & \path{run_cylinder3d_colagrossi.py} & $C_d$, $C_l$, surface profile vs Colagrossi 2018 \\
Oscillating cylinder (D\"utsch) & $\KC=5$, $\Re=100$, $U_{\max}=0.04$, domain $20D$ & \path{run_cylinder3d_oscillating.py} & Force history vs D\"utsch \\
Couette spin gate & $U_s=0.04$, $\Re_\Omega=100$, box $12D$ & \path{run_cylinder3d_spinning.py} & Analytic torque \\
Magnus & $\Re=100$, $U=0.05$, $30D\times16D$, spin ramped & \path{run_md_magnus_patch.py} & Lift vs literature anchor (to verify) \\
Sedimentation & $\rho_r=1.25$, $\Ga=19.6$ & \path{run_cylinder3d_sedimenting.py} & Terminal velocity vs reference \\
Vugts barge (H2-c) & $B/T=2$, forced heave/roll & \path{run_barge_roll.py} & Added mass / damping vs Vugts \\
TH01 hull & $L=256$, $\Fr$, $\Re_{\rm ref}$, tank G/Q/P zone presets & \path{run_hull_static.py}, \path{run_hull_speed.py} & Heave/trim/drag vs SYRF tank; BEM intercepts \\
Hydrofoil (Parkin) & $\Fr=0.95$, $h/c=1.8$, $c=128$, $\Re=4000$ & \path{run_hydrofoil_parkin.py} & Surface profile vs Parkin--Wu \\
\bottomrule
\end{longtable}

\paragraph{Convergence-ladder builder.} Refinement in this code means fixed-$\Cn$
($\xi\propto N$), fine-level mobility, and a steady-window extraction per rung. The
GUI generates the ladder from a card, estimates cost per rung, queues it, and
assembles the convergence plot and order estimate at the end
(\path{hull_ladder_analysis.py}, \path{plot_hull_ladder.py},
\path{make_hull_ladder_tex.py} are the existing pieces).

%======================================================================
\section{The trouble locator}\label{sec:locator}
%======================================================================
The hull driver already reports $\max|u-U|$ with its location relative to the CoG
and stops on a non-finite load. Make this a feature: on a NaN or a hot spot above
threshold, auto-checkpoint the last good state, stop, jump the camera to the cell,
drop a marker, and show the local $\phi$, $\tau$, safeguard flags, and whether the
cell sits in a band, a sponge, or a cut-link shell. This one feature will save more
user hours than any rendering option.

%======================================================================
\section{What sets it apart --- ranked}\label{sec:diff}
%======================================================================
\begin{enumerate}
\item \textbf{Recipe-first projects with lineage and envelope warnings.} No CFD GUI
  ships with its default state being a validated run and a gate.
\item \textbf{Ledgers and budgets as the health panel.} Mass ledger, energy-budget
  residual, two-channel force agreement, steady-window bracket, WB rest test.
\item \textbf{The feasibility triangle} ($N$, $\Ma$, $\tau$) with $\Cn$ and
  $\Pe_\phi$, live, with recipes plotted.
\item \textbf{The interface as a first-class object}: thickness histogram, spurious
  current, limiter and threshold activity maps. Scientific honesty built in.
\item \textbf{Strain rate is state}: dissipation, enstrophy, $Q$ exact; and the
  \emph{kinetic layer} ($S^{\rm neq}$, $a^{(3)}$, filtered $a^{(4)}$) as a
  ``where is the model working hardest'' map that only a moment-based Hermite solver
  can produce.
\item \textbf{The trouble locator.}
\item \textbf{The convergence-ladder builder} with cost per rung.
\item \textbf{One launcher for the Mac and the rack}, tenant-aware.
\item \textbf{Reproducibility by construction}: live CLI line, provenance
  \path{.npz}, \path{_viz} tag rule, figure sidecars.
\item \textbf{Four-layer pressure view} with the one-click WB rest test.
\end{enumerate}

%======================================================================
\section{Build sequence}\label{sec:build}
%======================================================================
\begin{description}[style=nextline]
\item[Phase 0 --- structure and units (no solver changes).] New sidebar; unit map
  with the three-form readouts; dimensionless-first setup form with the checks of
  \S\ref{sec:units}; backend badge; live CLI line; load existing provenance
  \path{.npz} archives and \path{kinetic_viz} snapshots. Ship the recipe cards as
  static project templates.
\item[Phase 1 --- monitors, forces, viewer (small driver additions).] Telemetry
  stream (\S\ref{sec:telemetry}) emitted by the drivers; Monitors page with the five
  health signals; forces with both channels and the bracket; viewer with the
  \path{kinetic_viz} layers and the air mask; timeline with four marker kinds;
  trouble locator; windowed MLUPS.
\item[Phase 2 --- validation and ladders.] Reference overlays and gates per card;
  convergence-ladder builder; cost model; envelope warnings from lineage.
\item[Phase 3 --- MD, queue, kinetic layer.] MD view; rack queue mirroring
  \path{status.sh}; safeguard activity maps; dissipation and divergence fields;
  energy-budget closure live; kinetic-layer fields; force-term layers.
\end{description}

%======================================================================
\section{Telemetry schema (so the GUI never imports \texorpdfstring{\code{torch}}{torch})}\label{sec:telemetry}
%======================================================================
Drivers print a JSON line every \path{--measure_every} steps (stdout or a
\path{.jsonl} beside the \path{.npz}); the GUI tails it. Proposed record; fields
marked $\dagger$ need a small driver addition, the rest are already computed.
\begin{lstlisting}
{"step": 12000, "t_lat": 12000, "t_phys": 0.84, "t_star": 2.31,
 "backend": "metal", "mlups_inst": 61.2, "mlups_cum": 63.0,
 "ma_inst": 0.071, "tau_min": 0.5083,
 "umax": 0.0412, "umax_loc": [212, 4, 96],
 "mass_ledger": {"phi": -1.3e-5, "levels": [-1.3e-5, 2.0e-6], "injected": 0.0},
 "ke": {"water": 3.1e-3, "air": 8.0e-6}, "pe_surface": 1.2e-3,
 "div_u_norm": 2.4e-4,
 "forces": {"Fx": -1.1e-2, "Fy": 0.0, "Fz": 0.42, "My": 3.0e-4,
            "Fx_p": -8.0e-3, "Fz_p": 0.41, "mea_Fx": -1.08e-2},
 "window": {"Fx": -1.10e-2, "Fx_prev": -1.12e-2},
 "safeguards": {"limiter_cells": 0, "threshold_cells": 0},
 "budget": {"W": 0.0, "D_near": 0.0, "D_far": 0.0, "D_air": 0.0,
            "Z_beach": 0.0, "Z_floor": 0.0, "dKE": 0.0, "dPE": 0.0, "res": 0.0},
 "snapshot": null, "checkpoint": null, "nonfinite": false}
\end{lstlisting}
$\dagger$: \path{mlups_inst}, \path{div_u_norm}, \code{safeguards},
\code{budget}, \path{pe_surface}, per-phase \code{ke}. Everything else is a
reformatting of what the drivers print today. A matching \path{run_spec.json}
(the settings the GUI serialises) should round-trip to the CLI line and be stored
inside the provenance \path{.npz} as a string.

%======================================================================
\section{Open questions for Colton and the frontend engineer}\label{sec:open}
%======================================================================
\begin{enumerate}
\item Process model: does the GUI launch the drivers as subprocesses (simplest;
  keeps the validated CLI path) or embed the solver? The telemetry schema assumes
  subprocesses and a tailed log; this also makes the rack a first-class target.
\item Rendering stack: port \path{kinetic_viz} (PyVista/VTK, Python) or
  reimplement the contract in the GUI's own renderer? If the GUI is web-based,
  the VTI export is the bridge.
\item Which drivers get the telemetry hook first? Suggest
  \path{run_hull_speed.py}, \path{run_cylinder3d_colagrossi.py},
  \path{run_rti3d_multiphase.py} --- they cover WB+FSI, MD, and single domain.
\item Where does the run spec live: inside the \path{.npz} only, or also a sidecar
  JSON beside it for grep-ability?
\item Energy-budget boxes (near/far/air) are case-specific in
  \path{h2_energy_budget.py}; the GUI needs a general box definition (body-relative
  radii) or a per-card default.
\item Unit map ownership: the drivers currently carry physical anchors ad hoc
  (\path{F_N}, \path{heave_mm}); a single \code{units} block in the run spec would
  let the GUI convert everything consistently.
\end{enumerate}

\appendix
%======================================================================
\section{Worked map: \code{run\_hull\_speed.py} flags \texorpdfstring{$\to$}{to} controls}
%======================================================================
\begin{longtable}{L{5.0cm}L{2.6cm}L{7.4cm}}
\caption{Every flag of \code{run\_hull\_speed.py} and where it lives in the
proposed architecture.}\\
\toprule
Flag & Page & Control \\
\midrule
\endfirsthead
\toprule
Flag & Page & Control \\
\midrule
\endhead
\path{--L}, \path{--config} & Geometry & Body length in cells; patch classification (CB / HKR / HKr / full) \\
\path{--Fn}, \path{--Re_ref}, \path{--g} & Fluids \& Interface & Dimensionless-first form; derived $U$, $\nu$, $\tau$ shown \\
\path{--xi}, \path{--Bo}, \path{--mobility}, \path{--rho_L} & Fluids \& Interface & Interface block; $\sigma$ derived from $\Bo$ \\
\path{--sink_cells}, \path{--trim}, \path{--no_equilibrate} & Bodies & Attitude offsets on top of the equilibrium solve \\
\path{--pad_up}, \path{--pad_down}, \path{--pad_side}, \path{--depth}, \path{--air} & Lattice & Tank layout in body lengths; drawn in the 3-D view \\
\path{--sponge}, \path{--xbeach}, \path{--xbeach_strength}, \path{--beach_y}, \path{--beach_gap}, \path{--zone_strength}, \path{--floor_fric} & Boundaries \& Zones & Zone regions with strength profiles; sized from $2\pi\Fr^2L$ \\
\path{--mass_correct}, \path{--pierce_bc} & Boundaries \& Zones & Mode selectors with one-line explanations \\
\path{--travel}, \path{--ramp_L}, \path{--no_frame_accel}, \path{--steps} & Run & Duration in body lengths; ramp; frame force toggle; step override \\
\path{--avg_L}, \path{--measure_every}, \path{--print_every} & Run / Monitors & Averaging window; telemetry cadence \\
\path{--band_cells}, \path{--refine}, \path{--sdf_backend}, \path{--cpt_path} & Geometry & SDF band and backend (libigl / CPT) \\
\path{--body_on_cpu}, \path{--no_gpu_kernels}, \path{--device} & Run & Backend and placement; fallback confirmation \\
\path{--init_state}, \path{--save_state} & Run / timeline & Warm start; restart checkpoint cadence \\
\path{--tag}, \path{--outdir}, \path{--smoke} & Run & Naming (the \path{_viz} rule), output folder, smoke preset \\
\bottomrule
\end{longtable}

\end{document}
